\hypertarget{appendix-11ba36693aa2}{}
\section*{Byte swapping and network order}
\section{Byte swapping and network order}\label{byte-swapping-and-network-order}

\textbf{Curriculum depth:} Introduced · \textbf{First seen:} C003 · \textbf{Latest development:} C003

\subsection{Definition and mechanism}\label{definition-and-mechanism}

Byte swapping reverses octet positions; host-to-network conversion instead adapts native order to a specified wire order. On an already-matching host conversion need not swap. Many classic Internet integer fields use big endian, but never infer every protocol's order from that convention. std::endian queries native order; C++23 std::byteswap reverses suitable integer bytes.

\subsection{Prerequisites}\label{prerequisites}

\begin{itemize}
\tightlist
\item
  \hyperref[appendix-bc07c3d25b18]{Endianness} --- Byte order tells how digit significance maps to sequence positions.
\item
  \hyperref[appendix-cc04044b5782]{Fixed-width integer types} --- The representation width fixes byte count and valid shift range.
\end{itemize}

\subsection{Used by}\label{used-by}

Update downstream uses when later chunks build on this concept.

\subsection{Recall and next depth}\label{recall-and-next-depth}

Explain why swapping twice restores x and why conversion may be a no-op.

\textbf{Pending:} Mixed/native scalar orders and protocol APIs.

\subsection{Source and navigation}\label{source-and-navigation}

\href{Continuous\%20Notes/Computing\%20Foundations\%20-\%20From\%20Bits\%20to\%20C++\%20and\%20Systems\%20APIs\#Where\%20the\%20bytes\%20live}{C003 teaching chapter} · \hyperref[appendix-fd090d982a15]{Part definitions} · \href{}{Dependency Map (02-Foundations-II, appendix-4ad47bdd269a)}
